\subsection{采用 ABE-backed 服务授权的理由}
\label{sec:abe-rationale-ch}

NDNSF 采用 ABE-backed 授权，是为了支持面向授权服务提供者的保密发现，
并将服务权限管理与身份签名凭证分离。这是针对 NDNSF 调用流程的设计取舍，
并非 ABE 全面优于角色授权的结论。执行授权决定实体是否可以调用或提供某项
服务；内容机密性决定谁能读取受保护的消息内容。签名验证、服务权限检查和
调用状态检查仍是不同的要求。

\paragraph{发现阶段的保密性。}
在选择 Provider 之前，User 发布 discovery descriptor（发现描述），其中
包含 Provider 决定是否返回 positive ACK 所需的位置、截止时间和资源要求。
例如，附近 UAV 需要知道观察区域和到达时限才能决定是否参与，但此时不需要
完整应用输入。ABE 无需逐一列出 Provider 身份，即可将发现描述的读取权限
限制为获准提供该服务的 Providers。这里的保密对象是缺少相应解密能力的外部
观察者；获准但最终未被选中的 Providers 仍可读取。内容加密不隐藏外层 Data
名称或流量模式。在请求级机密性路径中，完整应用输入及其解密材料通过
Selection 限定给选中的
Providers。Trust Schema 本身不加密内容；角色授权方案也可以配合 ABE 或
服务组加密来实现发现阶段的保密性。

\paragraph{权限聚合与凭证复用。}
KP-ABE 将属性关联到密文，将访问策略编码到解密密钥
（DKEY）中~\cite{dulal2023nacabe}。NDNSF 使用
\texttt{/SERVICE/<service>} 表示 Provider 的服务解密权限，使用
\texttt{/PERMISSION/<service>} 表示 User 的调用相关解密权限。在同一授权
机构和 ABE 参数代际内，一个实体的多项服务权限被合并为 OR 策略，编码到
单独为该实体签发的 DKEY 中，同时复用其身份签名凭证。DNMP 则通过角色签名
密钥和 Trust Schema 约束命令的目标与动作~\cite{dnmp}。已有角色可以覆盖多项
权限；不同权限需求可能要求调整角色分配、凭证或规则，但不意味着每项服务
都必须建立新角色。

作为对照，我们提出的每服务角色证书方案会为实体获准承担的每项服务角色签发
证书。User 使用相应服务角色密钥签名 Request 和 Selection Data，Provider
使用相应密钥签名 ACK 和 Response Data。验证方检查服务、消息类型、证书
所授角色、证书状态及调用状态，内容保密则由配套加密机制提供。ABE 避免了
这些额外的服务专用签名凭证，但 DKEY 仍需签发和更新，其大小及处理成本可能
随策略复杂度增长。证书方案也能复用凭证，因此凭证对象数量较少本身不足以
证明总成本更低。

\paragraph{为何继续使用 ABE-backed challenges？}
Requester Challenge（实现中的 \texttt{UserToken}）和 Provider Challenge
（\texttt{ProviderToken}）用于检查参与者在当前交换中是否具有相应的服务和
调用权限解密能力。签名标识参与身份，匹配检查和一次性状态检查拒绝不匹配或
重放的消息。这沿用了属性化权限管理，而不必另外增加服务角色签名凭证，但
并非没有额外成本：保密发现已需要 Provider DKEY，而 User 的调用权限 DKEY
和 challenge 处理仍是额外授权成本。ABE 保密配合角色证书授权同样是有效的
替代方案。

\subsection{运行时权限撤销}
\label{sec:authorization-withdrawal-ch}

当前运行时通过 Controller 签名的授权状态以及调用处理点上的检查实现权限
撤销。参与者需要获取并验证更新后的状态，才能执行新公布的权限撤销。
ABE 参数轮换和 DKEY 更新保护随后使用新代际加密的交换。目前轮换作用于整个
Controller 的 ABE 参数代际，其他仍获授权的成员也可能需要刷新 DKEY；
这不是无需更新其他成员的撤销方案。

DNMP 描述了通过轮换加密密钥限制后续测量数据访问，但没有详细给出可对应
比较的在线命令授权撤销协议~\cite{dnmp}。角色证书对照方案可以撤销相关凭证，
并要求执行节点检查签名的撤销状态或拒绝过期证书。这是在验证方取消凭证
所代表的权限，而不是物理取回已分发的证书。两类方案都需要传播更新信息，
也都无法收回已披露的密钥或明文。

定向集成测试已检验 NDNSF 的运行时撤销检查，包括身份级、证书级和服务级
撤销。这些测试支持已覆盖运行时边界上的行为，不代表整个网络中的瞬时撤销，
也不能证明其成本低于证书方案。后续评估应测量状态传播延迟、拒绝生效延迟、
重新分发密钥的流量，以及对未撤销成员的影响。因此，运行时状态验证和密钥
轮换使 NDNSF 支持权限撤销，同时带来明确的状态传播和重新分发密钥的成本。
